\subsection*{Problem 1.2}

\subsubsection*{(a)}
Each announcement by Tandon becomes common knowledge. 
Consider the repeated announcement: 
\emph{``Does any of you know whether you have 
mud on your own forehead?''}. 
If no child says ``yes'' after the first repetition, 
then all children know that there are 
\textbf{more than one} muddy child. 
The justification is simple: each child sees the others' 
foreheads. 
If there were exactly one muddy child, that 
child would see all others clean and would therefore know 
his own muddiness immediately after the first announcement 
\emph{``At least one of you has mud on your forehead''}. 
Since he does not say ``yes'', the situation 
with a single muddy child is impossible.

Similarly, if no child says ``yes'' after the second 
repetition, all children deduce there are 
\textbf{more than two} muddy children. By induction, 
after the $i$-th repetition with no ``yes'', 
the children know there must be \textbf{more than $i$} 
muddy children. Conversely, if a child knows 
there are more than $i-1$ muddy children 
(from the previous round's silence) and sees only $i-2$ muddy 
children among the others, he deduces he himself must be muddy. 
Hence, after $k$ repetitions, 
every child knows whether he has mud on his forehead,
because the set of possible worlds will become empty.

\paragraph*{Initial Kripke structure:}
Define \( M = \langle S, \sim, V \rangle \) where
\begin{gather*}
    S = \{0,1\}^n, \\
    \text{Agents} = \{c_i \mid 1 \leq i \leq n\}, \\
    \forall i \in \{1,\dots,n\}: 
        \sim_{c_i}(l1r) = \{l1r, l0r\} = \sim_{c_i}(l0r)
        \quad \text{ and } \lr{i}{0,1}\\
    \forall i \in \{1,\dots,n\}: 
        V(0_{c_i}) = \{ l0r : \lr{i}{0,1} \},\quad 
        V(1_{c_i}) = S \setminus V(0_{c_i}).
        \tag{1.2(a1)}\label{eq:1.2(a1)}
\end{gather*}
The Kripke structure shows all possible worlds as 
$n$-length bit strings, where the $i$-th bit 0 means child 
$c_i$ is clean, otherwise muddy. The equivalence relation $\sim$ 
reflects the fact that each child can see the others' foreheads 
but not their own. The proposition $0_{c_i}$ represents that 
child $c_i$ is clean, and $V(0_{c_i})$ is the set of worlds 
where that holds.

Now we examine how each announcement by Tandon restricts the 
Kripke structure $M$. Let $\phi_0$ denote the announcement 
\emph{``At least one of you has mud on your forehead''}, 
and let $\phi$ denote the announcement 
\emph{``Does any of you know whether you have mud on your own forehead?''}. 
We write $\phi_i$ (with $i \ge 1$) for the $i$-th repetition of $\phi$.

\paragraph*{After the announcement $\phi_0$:}
We encode $\phi_0$ as $\phi_0 = \an$. The restricted Kripke structure 
is $M^{(0)} = M|_{\phi_0}$ where
\begin{gather*}
    S^{(0)} \coloneqq [\![ \phi_0 ]\!] \coloneqq 
        \{ s \in S : M,s \models \phi_0 \}, \\
    \forall i \in \{1,\dots,n\}: 
        \sim_{c_i}^{(0)} \coloneqq 
        \sim_{c_i} \cap ([\![ \phi_0 ]\!] \times [\![ \phi_0 ]\!]), \\
    V^{(0)}(p) \coloneqq V(p) \cap [\![ \phi_0 ]\!].
    \tag{1.2(a2)}\label{eq:1.2(a2)}
\end{gather*}
For any world $s \in S$,
\begin{align*}
    M,s \models \phi_0 
        &\iff M,s \models \an \\
        &\iff \bigvee_{i=1}^{n} M,s \models 1_{c_i} \\
        &\iff \bigvee_{i=1}^{n} s \in V(1_{c_i}). \tag{1.2(a3)}\label{eq:1.2(a3)}
\end{align*}
Except the world $\{0\}^n$, all worlds satisfy \eq{1.2(a3)}; hence only $\{0\}^n$ 
is eliminated from $M^{(0)}$. Thus $S^{(0)} = S \setminus \{0\}^n$, and 
the relations and valuation are restricted accordingly.

\paragraph{Intuitive verification:} The announcement itself forces 
the condition that there is at least one muddy child, so the world 
with no muddy child becomes impossible. All other worlds have at least 
one muddy child and therefore survive.

\paragraph*{After the announcement $\phi_i$ for $i \ge 1$:}
By construction $\phi_i = \phi$ for all $i \ge 1$; the index $i$ denotes 
the repetition round. We encode $\phi$ as $\phi = \ani{j}$. This formula 
only allows worlds where no child knows whether he is muddy.

After the $i$-th repetition, the restricted Kripke structure is 
$M^{(i)} = M^{(i-1)}|_{\phi_i} = \langle S^{(i)}, \sim^{(i)}, V^{(i)} \rangle$ with
\begin{gather*}
    S^{(i)} \coloneqq [\![ \phi_i ]\!] \coloneqq 
        \{ s \in S^{(i-1)} : M^{(i-1)},s \models \phi_i \}, \\
    \forall i \in \{1,\dots,n\}: 
        \sim_{c_i}^{(i)} \coloneqq 
        \sim_{c_i}^{(i-1)} \cap ([\![ \phi_i ]\!] \times [\![ \phi_i ]\!]), \\
    V^{(i)}(p) \coloneqq V^{(i-1)}(p) \cap [\![ \phi_i ]\!].
    \tag{1.2(a4)}\label{eq:1.2(a4)}
\end{gather*}
For any world $s \in S^{(i-1)}$ and $i = 1,2,\dots$,
\begin{align*}
    M^{(i-1)},s \models \phi_i 
        &\iff M^{(i-1)},s \models \ani{j} \\
        &\iff \bigwedge_{j=1}^{n} \big( 
            M^{(i-1)},s \models \neg K_{c_j}1_{c_j} 
            \;\land\; M^{(i-1)},s \models \neg K_{c_j}0_{c_j} \big) \\
        &\iff \bigwedge_{j=1}^{n} \big( 
            \neg (M^{(i-1)},s \models K_{c_j}1_{c_j}) 
            \;\land\; \neg (M^{(i-1)},s \models K_{c_j}0_{c_j}) \big) \\
        \iff \bigwedge_{j=1}^{n} \big( 
            &\neg(\forall t \in \sim_{c_j}^{(i-1)}(s): t \in V^{(i-1)}(1_{c_j}))
            \;\land\; \\
            &\neg(\forall t \in \sim_{c_j}^{(i-1)}(s): 
                t \in V^{(i-1)}(0_{c_j})) \big). \tag{1.2(a5)}\label{eq:1.2(a5)}
\end{align*}

\subparagraph{Case $i=1$ (first repetition):} 
Let $m_{j1} = \{ l1r : \lr{j}{0} \}$ be the set of
 worlds where only child $c_j$ is muddy, 
and $m_1 = \bigcup_{j=1}^{n} m_{j1}$ the set 
of worlds with exactly one muddy child. 
Observe that \textbf{for all $j$},
\begin{align*}
    &\forall l1r \in m_{j1}: \quad 
        \sim_{c_j}^{(0)}(l1r) = \{l1r\} \quad 
            (\text{since } l0r \notin S^{(0)}), \tag{1.2(a6)}\label{eq:1.2(a6)}\\
    &\forall l1r,l0r \in S^{(0)} \setminus m_1 \text{ and, } \lr{j}{0,1}: \\ 
        &\sim_{c_j}^{(0)}(l1r) = \{l1r,l0r\} = \sim_{c_j}^{(0)}(l0r). 
        \tag{1.2(a7)}\label{eq:1.2(a7)}
\end{align*}
From \eq{1.2(a5)} and \eq{1.2(a6)}, for any child $c_j$ and any $s \in m_{j1}$,
we have
$M^{(0)},s \models K_{c_j}1_{c_j} (\because l1r \in 1_{c_j})$
 and $M^{(0)},s 
\not\models K_{c_j}0_{c_j} (\because l1r \not\in V^{(0)}(0_{c_j})) $. 
Hence all worlds in $m_{j1}$ fail to satisfy 
$\phi_1$ and are eliminated from $M^{(1)}$. 
This elimination is possible only because the world 
$\{0\}^n$ was already removed from $S^{(0)}$, 
which makes the equivalence classes of worlds in $m_{j1}$ singletons. 
For any world $s \in S^{(0)} \setminus m_1$, 
from \eq{1.2(a7)}
\( (l1r \not\in V^{(0)}(0_{c_j})\text{ and }l0r \not\in V^{(0)}(1_{c_j}) )\),
where $"l1r"$ and $"l0r"$ belongs to $\sim_{c_j}^{(0)}(s)$.
Hence, we obtain
\[
\bigwedge_{j=1}^{n} \big( M^{(0)},s \not\models K_{c_j}1_{c_j} 
\;\land\; M^{(0)},s \not\models K_{c_j}0_{c_j} \big),
\]

so they satisfy $\phi_1$ and survive. Consequently,
$S^{(1)} = S^{(0)} \setminus m_1$, 
i.e., all worlds with exactly one muddy 
child are removed after the first repetition.

\subparagraph{Case $i=2$ (second repetition):} 
By the same reasoning as in case $i=1$, 
because worlds with exactly one muddy child were eliminated 
after $\phi_1$, the equivalence classes for worlds with exactly 
two muddy children become singletons. 
Therefore, after the second repetition $\phi_2$, 
all worlds with exactly two muddy children are eliminated.

We can frame this as a proof by induction: 
after the $i$-th repetition of $\phi$, 
only worlds with \textbf{more than $i$ muddy children} 
remain in $M^{(i)}$. Hence, after $k$ repetitions, 
all children know their own muddy status,
because the set of possible worlds will become empty.

\paragraph{Intuitive verification:} If there were exactly one muddy child, 
then after the announcement $\phi_0$, that child would know 
he is muddy because he sees all others clean. Similarly,
if no child says ``yes'' after the first repetition, 
then any muddy child must see at least one other muddy child; 
otherwise he would have deduced his own muddiness. 
This reasoning repeats: after the $i$-th repetition with no 
``yes'', each muddy child sees at least $i$ other muddy children, 
so he still cannot be certain. But once the number of 
repetitions reaches the actual number of muddy children, 
every muddy child sees only $k-1$ muddy others 
and deduces he must be muddy. Hence, after $k$ repetitions, 
all $k$ muddy children know they are muddy,
because the set of possible worlds will become empty.

\subsubsection*{(b)}

\paragraph*{Initial Kripke structure:}
Define \( M = \langle S, \sim, V \rangle \) with the same set \(S\) and 
valuations as in (a), but now all children have bandages over their eyes, 
so they cannot distinguish any worlds:
\begin{gather*}
    S = \{0,1\}^n, \\
    \text{Agents} = \{c_i \mid 1 \leq i \leq n\}, \\
    \forall i \in \{1,\dots,n\}: \sim_{c_i} = S \times S, \\
    \forall i \in \{1,\dots,n\}: 
        V(0_{c_i}) = \{ l0r : \lr{i}{0,1} \},\quad 
        V(1_{c_i}) = S \setminus V(0_{c_i}).
\end{gather*}

\paragraph*{After announcement $\phi_0$:}
As shown in part (a) \eq{1.2(a3)}, after the announcement $\phi_0$, 
the restricted Kripke structure \( M^{(0)} = M|_{\phi_0} \) satisfies
\begin{gather*}
    S^{(0)} = S \setminus \{0\}^n, \\
    \forall i \in \{1,\dots,n\}:
    \sim_{c_i}^{(0)} = S^{(0)} \times S^{(0)}. 
        \tag{1.2(b1)}\label{eq:1.2(b1)}
\end{gather*}
The valuation function is restricted accordingly, following the definition 
from part (a) \eq{1.2(a2)}.

\paragraph*{After announcement $\phi_i$ for $i \ge 1$:}
After the $i$-th repetition, the restricted Kripke structure is 
\( M^{(i)} = M^{(i-1)}|_{\phi_i} = 
    \langle S^{(i)}, \sim^{(i)}, V^{(i)} \rangle \), defined 
similarly to part (a) \eq{1.2(a4)}.

\subparagraph{Case $i=1$:} 
Since the children cannot see, using formula \eq{1.2(a4)} and the 
equivalence relation \eq{1.2(b1)}, for any world \( s \in S^{(0)} \):
\begin{align*}
    &\forall j \in \{1,\dots,n\}: 
        \exists t_1, t_2 \in \sim_{c_j}^{(0)}(s): 
        t_1 \notin V^{(0)}(1_{c_j}) \;\land\; t_2 \notin V^{(0)}(0_{c_j}) \\
        &\because (t_1 \in S^{(0)} \setminus V^{(0)}(1_{c_j}) \text{ and }
         t_2 \in S^{(0)} \setminus V^{(0)}(0_{c_j})) \\
    \iff &\forall j \in \{1,\dots,n\}: 
        \neg (M^{(0)}, s \models K_{c_j}1_{c_j}) 
        \;\land\; \neg (M^{(0)}, s \models K_{c_j}0_{c_j}) \\
    \iff &\forall j \in \{1,\dots,n\}: 
        M^{(0)}, s \models \neg K_{c_j}1_{c_j} \land \neg K_{c_j}0_{c_j} \\
    \iff &\bigwedge_{j=1}^{n} M^{(0)}, s \models \neg K_{c_j}1_{c_j} 
        \land \neg K_{c_j}0_{c_j}.
\end{align*}
Thus every world in \( S^{(0)} \) satisfies the announcement \(\phi_1\), so 
\( S^{(1)} \coloneqq [\![ \phi_1 ]\!] = S^{(0)} \). This reflects the fact 
that in every world of \( S^{(0)} \), no child knows his muddy status. 
In fact, \( M^{(1)} = M^{(0)} \).

\subparagraph{Case $i=2,\dots$:} 
The same pattern holds for all further repetitions of \(\phi\); 
(because the structure $M^{(i-1)}$ and announcement $\phi$ is
same for all future repetition, so we'll get same result)
the Kripke 
structure remains unchanged after each round. Hence,
\[
\forall i \in \{2,\dots\}: M^{(i)} = M^{(i-1)} = M^{(0)}.
\]

\paragraph{Intuitive verification:} Since no child can see the status of 
others, they lack the necessary information to determine their own 
muddiness. Their only knowledge about the actual world is that there 
is at least one muddy child (from \(\phi_0\)). Consequently, they will 
never learn their own muddy status, even after infinitely many repetitions 
of \(\phi\).

\subsubsection*{(c)}

\paragraph*{Initial Kripke structure:} 
We use the same initial Kripke structure as in $M$ defined in \eq{1.2(a1)}.
To encode Tandon's statement, we define an action model which
captures the inattentiveness of \emph{child 1}. 
Define \( M^{a}= \langle S^{a},\sim^{a},prec \rangle \) such that
\begin{gather*}
    S^{a} = \{A,B\}, \\
    \sim^{a} = \{(s,s) : s \in S^{a}\}, \\
    prec(A) = \bigvee_{j=1}^{n} 1_{c_j}, \\
    prec(B) = \text{True}.
\end{gather*}
In the action model $M^{a}$, $A$ represents the announcement 
\emph{``At least one of you has mud on your forehead''} and 
$B$ represents a vacuous statement (always true).

\paragraph*{After Tandon's statement:}
The modified Kripke structure after Tandon's statement
is the product update of $M$ and $M^a$:
$M'^{(0)} = M \otimes M^a = \langle S'^{(0)},\sim'^{(0)},V'^{(0)} \rangle$,
where
\begin{gather*}
    S'^{(0)} = \{ (s,a) : s \in S,\; a \in S^a,\; M,s \models prec(a) \}, \\
    \text{Agents} = \{c_j : j \in \{1,2,\dots,n\}\}, \\
    \sim'^{(0)}_{c_1} = \{ ((s_1,a_1),(s_2,a_2)) : 
        \forall s_1 \in S,\; \forall s_2 \in \sim_{c_1}(s_1),\;
        \forall a_1,a_2 \in S^a \}, \tag{1.2(c1)}\label{eq:1.2(c1)}\\
    \forall j \in \{2,3,\dots,n\}:\;
        \sim'^{(0)}_{c_j} = \{ ((s_1,a),(s_2,a)) :
            \forall s_1 \in S,\; \forall s_2 \in \sim_{c_1}(s_1),\;
            \forall a \in S^a \}, \tag{1.2(c2)}\label{eq:1.2(c2)}\\
    V'^{(0)}(p) = \{ (s,a) : \forall s \in V(p),\; \forall a \in S^a \}.
\end{gather*}
The equivalence relation for child $c_1$ differs from 
that of other children because $c_1$ cannot distinguish 
between $A$ and $B$, whereas the others can.
The restricted structure does not contain the world $(\{0\}^n,A)$ 
because, by construction of the product update, 
$M, \{0\}^n \not\models prec(A)$, while $prec(B)=\text{True}$; 
thus $(s,B) \in S'$ for all $s \in S$.

\paragraph*{(i)}
For $n=2$, the Kripke structure is as follows.

\textbf{Before Tandon's statement $M$:}
\begin{align*}
    &S = \{ 00, 01, 10, 11 \}, \\
    &\sim_{c_1}(00) = \{00,10\} = \sim_{c_1}(10), \\
    &\sim_{c_1}(01) = \{01,11\} = \sim_{c_1}(11), \\
    &\sim_{c_2}(00) = \{00,01\} = \sim_{c_2}(01), \\
    &\sim_{c_2}(01) = \{01,11\} = \sim_{c_2}(11), \\
    &V(0_{c_1}) = \{ 00, 01 \},\quad V(1_{c_1}) = S \setminus V(0_{c_1}), \\
    &V(0_{c_2}) = \{ 00, 10 \},\quad V(1_{c_2}) = S \setminus V(0_{c_2}).
\end{align*}

\textbf{After Tandon's statement $M'^{(0)}$:} 
The product update $M \otimes M^a$ yields
\begin{align*}
    &S'^{(0)} = \{ (01,A), (10,A), (11,A),
           (00,B), (01,B), (10,B), (11,B) \}, \\
    &\forall a \in \{A,B\}: \\
    &\quad \sim'^{(0)}_{c_1}(01,a) = \{(01,A),(11,A),(01,B),(11,B)\} 
        = \sim'^{(0)}_{c_1}(11,a), \\
    &\quad \sim'^{(0)}_{c_1}(10,a) = \{(10,A),(00,B),(10,B)\} 
        = \sim'^{(0)}_{c_1}(00,B), \\
    &\quad \sim'^{(0)}_{c_2}(10,a) = \{(10,a),(11,a)\} = \sim'^{(0)}_{c_2}(11,a), \\
    &\quad \sim'^{(0)}_{c_2}(00,B) = \{(00,B),(01,B)\} = \sim'^{(0)}_{c_2}(01,B), \\
    &\quad \sim'^{(0)}_{c_2}(01,A) = \{(01,A)\}, \\
    &V'^{(0)}(0_{c_1}) = \{ (01,A), (00,B), (01,B) \},\quad
        V'^{(0)}(1_{c_1}) = S'^{(0)} \setminus V'^{(0)}(0_{c_1}), \\
    &V'^{(0)}(0_{c_2}) = \{ (00,B), (10,A), (10,B) \},\quad
        V'^{(0)}(1_{c_2}) = S'^{(0)} \setminus V'^{(0)}(0_{c_2}).
\end{align*}

\paragraph*{(ii)}

\textbf{Is it possible for the children to ascertain whether they are muddy?}

To answer this, we analyse what happens after each repetition round of 
the second announcement $\phi$, 
\emph{``Does any of you know whether you have mud on your own forehead?''}.
After the $i$-th repetition ($i \ge 1$) of $\phi$ (denoted $\phi_i$), 
the restricted Kripke structure is 
$M'^{(i)} = M'^{(i-1)}|_{\phi_i} = \langle S'^{(i)}, 
\sim^{(i)}, V'^{(i)} \rangle$ with
\begin{gather*}
    S'^{(i)} \coloneqq [\![ \phi_i ]\!] \coloneqq 
    \{ s \in S'^{(i-1)} : M'^{(i-1)},s \models \phi_i \}, \\
    \forall i \in \{1,\dots,n\}: 
    \sim_{c_i}^{(i)} \coloneqq 
    \sim_{c_i}^{(i-1)} \cap ([\![ \phi_i ]\!] \times [\![ \phi_i ]\!]), \\
    V'^{(i)}(p) \coloneqq V'^{(i-1)}(p) \cap [\![ \phi_i ]\!].
\end{gather*}
The encoding of $\phi$ and the formula are similar to \eq{1.2(a4)} and \eq{1.2(a5)}:
\begin{align*}
    M'^{(i-1)},s \models \phi_i 
    &\iff M'^{(i-1)},s \models \ani{j} \\
    \iff \bigwedge_{j=1}^{n} \Big( 
        &\neg\bigl(\forall t \in \sim_{c_j}'^{(i-1)}(s): 
            t \in V'^{(i-1)}(1_{c_j})\bigr)
        \;\land\; \\
        &\neg\bigl(\forall t \in \sim_{c_j}'^{(i-1)}(s): 
            t \in V'^{(i-1)}(0_{c_j})\bigr) \Big). 
        \tag{1.2(c3)}\label{eq:1.2(c3)}
\end{align*}

\subparagraph{Case $i=1$ (first repetition):} 
Define 
\(m_{ja1}' = \{ (l1r,A) : \lr{j}{0} \}\) 
(worlds in $S'$ with only child $j$ muddy)
and
\(m_{a1}' = \bigcup_{j=2}^{n} m_{ja1}'\)
(the set of worlds where exactly one child other than $c_1$ is muddy).

Observe that for $j=1$,
\begin{align*}
    &\forall (l1r,A) \in m_{1a1}': \\ 
        &\sim_{c_1}'^{(0)}((l1r,A)) = \{(l1r,A),(l0r,B),(l1r,B)\} 
        = \sim_{c_1}'^{(0)}((l0r,B)), \tag{1.2(c4)}\label{eq:1.2(c4)}\\
    &\forall (l1r,a),(l0r,b) \in S'^{(0)} \setminus m'_{a1} 
        \text{ with } \lr{j}{0,1}: \\ 
        &\sim_{c_1}'^{(0)}((l1r,a)) = 
        \{(l1r,a),(l1r,b),(l0r,a),(l0r,b)\} = \sim_{c_j}'^{(0)}((l0r,b)). 
        \tag{1.2(c5)}\label{eq:1.2(c5)}
\end{align*}
For all $j \ge 2$,
\begin{align*}
    \forall j \in &\{2,3,\dots,n\}: \\
    &\forall (l1r,A) \in m_{ja1}': \quad 
        \sim_{c_j}'^{(0)}((l1r,A)) = \{(l1r,A)\} 
        \quad (\text{since } (l0r,A) \notin S'^{(0)}), 
            \tag{1.2(c6)}\label{eq:1.2(c6)}\\
    &\forall (l1r,a),(l0r,a) \in S'^{(0)} \setminus m'_{a1} 
        \text{ with } \lr{j}{0,1}: \\ 
        &\sim_{c_j}'^{(0)}((l1r,a)) = \{(l1r,a),(l0r,a)\} = 
        \sim_{c_j}'^{(0)}((l0r,a)). \tag{1.2(c7)}\label{eq:1.2(c7)}
\end{align*}

From \eq{1.2(c3)} and \eq{1.2(c6)}, 
for any child $c_j$ ($j \ge 2$) and any $s \in m_{ja1}'$,
we have $M'^{(0)},s \models K_{c_j}1_{c_j}$ 
(because $(l1r,A) \notin V'^{(0)}(1_{c_j})$) and 
$M'^{(0)},s \not\models K_{c_j}0_{c_j}$ 
(because $(l1r,A) \notin V'^{(0)}(0_{c_j})$). 
Hence all worlds in $m_{ja1}'$ fail to satisfy $\phi_1(=\phi)$ 
and are eliminated from $M'^{(1)}$. This elimination is possible 
only because the world $(\{0\}^n,A)$ was already removed from $S^{(0)}$, 
which makes the equivalence classes of worlds in $m_{ja1}'$ singletons.

For any world $s \in S^{(0)} \setminus m_{a1}'$, using \eq{1.2(c4)}, 
\eq{1.2(c5)} and \eq{1.2(c7)}, we can argue that we will always 
find worlds $t_1,t_2 \in \sim_{c_j}'^{(0)}(s)$ for every $j \ge 1$ 
such that $t_1 \notin V'^{(0)}(1_{c_j})$ and $t_2 \notin V'^{(0)}(0_{c_j})$. 
Hence
\[
\bigwedge_{j=1}^{n} \big( M'^{(0)},s \not\models K_{c_j}1_{c_j} 
\;\land\; M'^{(0)},s \not\models K_{c_j}0_{c_j} \big),
\]
so they satisfy $\phi_1$ and survive. Consequently,
$S^{(1)} = S^{(0)} \setminus m_{a1}'$; i.e., all \textbf{$A$-worlds} 
(worlds of the form $(s,A)$) with \textbf{exactly one} muddy child 
(other than $c_1$) are removed after the first repetition, 
while $A$-worlds with two or more muddy children and all $B$-worlds survive.

\subparagraph{Case $i=2$ (second repetition):} 
By the same reasoning as in case $i=1$, because $A$-worlds 
with exactly one muddy child were eliminated after $\phi_1$, t
he equivalence classes for worlds with exactly two muddy children 
become singletons. Therefore, after the second repetition $\phi_2$, 
all $A$-worlds with exactly two muddy children (other than $c_1$) 
are eliminated. Again, $A$-worlds where $c_1$ is muddy and 
all $B$-worlds survive.

We can frame this as a proof by induction: after the $i$-th repetition of 
$\phi$, the surviving worlds are those $A$-worlds with 
\textbf{more than $i$ muddy children} or with \textbf{child $c_1$ muddy}, 
together with \textbf{all $B$-worlds}. Hence, even after infinitely 
many repetitions of $\phi$, $A$-worlds with child $c_1$ muddy and 
all $B$-worlds remain in the restricted Kripke structure. 
Other children can ascertain their muddy status only if 
Tandon announced $A$; otherwise, no one can deduce. 
Child $c_1$ can never deduce whether he has mud on his forehead in any case.

\paragraph{Intuitive verification:} 
Since all children except child $c_1$ are attentive, 
they know for sure which announcement Tandon made: $A$ 
(``At least one of you has mud on his forehead'') or 
$B$ (vacuous statement). Therefore, with each repetition 
they can deduce whether they are muddy (as in the previous cases). 
However, child $c_1$ is inattentive and never knows which 
announcement was made, 
so he can never deduce his own muddy status. 
In the case of a vacuous statement $B$, there is insufficient 
information about the true world, for any child to ever deduces anything.

\subsubsection*{(d)}

We assume the same initial Kripke structure as in \eq{1.2(a1)}. 
The only change is that the initial announcement $\phi_0$ was 
\emph{``At least one of you has mud on your forehead''}, but now it is 
changed to \emph{``Child 1 has mud on his forehead''}. 
Accordingly, the encoding of the announcement becomes $\phi_0 = 1_{c_1}$.

\paragraph*{After announcement $\phi_0 = 1_{c_1}$:}
The restricted Kripke structure $M^{(0)} = M|_{\phi_0}$ is defined as in 
\eq{1.2(a2)} with $S^{(0)} \coloneqq [\![ \phi_0 ]\!] = 
\{ s \in S : M,s \models \phi_0 \}$. 
For any world $s \in S$,
\begin{align*}
    M,s \models \phi_0 
        &\iff M,s \models 1_{c_1} \\
        &\iff s \in V(1_{c_1}). \tag{1.2(d1)}\label{eq:1.2(d1)}
\end{align*}
From \eq{1.2(d1)} only the worlds in $V(1_{c_1})$ survive the announcement, hence
\begin{align*}
    &S^{(0)} \coloneqq [\![ \phi_0 ]\!] = V(1_{c_1}), \\
    &\sim_{c_1}^{(0)}(1r) = \{1r\} \quad \text{for all } r \in \{0,1\}^{n-1}, \\
    &\forall j \in \{2,\dots,n\}: \\
    &\qquad \sim_{c_j}^{(0)}(1l1r) = \{ 1l1r, 1l0r \} = \sim_{c_j}^{(0)}(1l0r) \\
    &\qquad \text{where } \forall l \in \{0,1\}^{j-1},\; 
        \forall r \in \{0,1\}^{n-j-1}. \tag{1.2(d2)}\label{eq:1.2(d2)}
\end{align*}
The valuation function is restricted similarly.

\paragraph*{Can others deduce whether they are muddy?}
To answer this, we check which worlds allow at least one child other than 
\emph{child 1} to deduce his muddy status. For that, we restrict the Kripke structure 
against the repetitive announcement 
\[
\phi = \bigwedge_{j=2}^{n} \big( \neg K_{c_j}1_{c_j} \land \neg K_{c_j}0_{c_j} \big).
\]
After the $i$-th repetition, the restricted structure is 
$M^{(i)} = M^{(i-1)}|_{\phi_i} = \langle S^{(i)}, \sim^{(i)}, V^{(i)} \rangle$. 
Worlds in which any other child knows his own state are eliminated by $\phi$. 
Since the announcement is similar to that in \eq{1.2(a5)}, the formula becomes:
\begin{align*}
    M^{(i-1)},s \models \phi_i 
        &\iff M^{(i-1)},s \models \bigwedge_{j=2}^{n} 
            \big( \neg K_{c_j}1_{c_j} \land \neg K_{c_j}0_{c_j} \big) \\
        &\iff \bigwedge_{j=2}^{n} \Big( 
            \neg\bigl(\forall t \in \sim_{c_j}^{(i-1)}(s): 
                t \in V^{(i-1)}(1_{c_j})\bigr) \\
        &\qquad\quad \land\; 
            \neg\bigl(\forall t \in \sim_{c_j}^{(i-1)}(s): 
                t \in V^{(i-1)}(0_{c_j})\bigr) \Big). 
                \tag{1.2(d3)}\label{eq:1.2(d3)}
\end{align*}

\subparagraph{Case $i=1$ (first repetition):} 
Analysing the equivalence relation $\sim^{(0)}$ 
from \eq{1.2(d2)}, for any $s \in S^{(0)}$ we have:
\begin{align*}
    &\forall j \in \{2,\dots,n\}: 
        \exists t_1, t_2 \in \sim_{c_j}^{(0)}(s):
            t_1 \notin V^{(0)}(1_{c_j}) \;\land\; t_2 \notin V^{(0)}(0_{c_j}) \\
    &\qquad (\text{choose } t_1 = 1l0r,\; t_2 = 1l1r \text{ with appropriate } l,r) \\
    \iff &\forall j \in \{2,\dots,n\}: 
        \neg\bigl(\forall t \in \sim_{c_j}^{(0)}(s): 
            t \in V^{(0)}(1_{c_j})\bigr)
        \;\land\; 
        \neg\bigl(\forall t \in \sim_{c_j}^{(0)}(s): 
            t \in V^{(0)}(0_{c_j})\bigr) \\
    \iff &\bigwedge_{j=2}^{n} 
        \Big( \neg (M^{(0)},s \models K_{c_j}1_{c_j})
        \;\land\; \neg (M^{(0)},s \models K_{c_j}0_{c_j}) \Big) \\
    \iff &\bigwedge_{j=2}^{n} 
        \big( M^{(0)},s \models \neg K_{c_j}1_{c_j} \land \neg K_{c_j}0_{c_j} \big). 
        \tag{1.2(d4)}\label{eq:1.2(d4)}
\end{align*}
Since every world $s \in S^{(0)}$ satisfies \eq{1.2(d4)}, 
all worlds survive the announcement. 
Thus $S^{(1)} \coloneqq [\![ \phi_1 ]\!] = S^{(0)}$ and 
$M^{(1)} = M^{(0)}|_{\phi_1} = M^{(0)}$. 
This reflects that, even after the first repetition, 
no child other than \emph{child 1} can deduce his muddy status.

\subparagraph{Case $i=2,\dots$:} 
The same pattern holds for all further repetitions of 
$\phi$ because the structure $M^{(i-1)}$ and the announcement 
$\phi$ remain identical each round. Hence,
\[
\forall i \ge 2:\quad M^{(i)} = M^{(i-1)} = M^{(0)}.
\]
This shows mathematically that even after infinitely 
many repetitions of the announcement $\phi$, 
no child other than \emph{child 1} can deduce whether he is muddy.

\paragraph{Intuitive verification:} The announcement 
\emph{``Child 1 has mud on his forehead''} ($\phi_{c1}$) 
initially conveys more information 
than the announcement \emph{``At least one child has mud''} ($\phi_0$). 
However, as repetitions of the second announcement ($\phi$) proceed, 
$\phi_{c1}$ does not propagate any further information, 
whereas $\phi_0$ propagates the information that 
\textbf{more than $i-1$} children are muddy at round $i$. 
Since $\phi_0$ unfolds additional information with each repetition, 
the children were able to deduce their own muddy status, 
which was not possible with $\phi_{c1}$.